\documentclass[10pt,a4paper]{amsart}
\usepackage[margin=19mm]{geometry}
\usepackage{amsmath,amssymb,mathtools}
\usepackage[scr=rsfs,cal=euler]{mathalfa}
\usepackage{xfrac}
\usepackage[unicode,hidelinks,bookmarks=false]{hyperref}
\newcommand{\ie}{i.e.\ }
\newcommand{\cf}{cf.\ }
\newcommand{\resp}{respectively\ }
\newcommand{\Subsection}{\subsection}
\providecommand{\nobreakdash}{\nobreak\hbox{-}}
\setlength{\emergencystretch}{2em}
\multlinegap0pt
\usepackage{amssymb}
\usepackage{mathtools}
\usepackage[shortlabels]{enumitem}
\newtheorem{proposition}{Proposition}
\newtheorem{theorem}{Theorem}
\usepackage{cleveref}
\crefname{proposition}{Proposition}{Propositions}
\crefname{theorem}{Theorem}{Theorems}
\DeclareMathOperator{\Law}{Law}
\newcommand{\R}{\mathbb R}
\newcommand{\cH}{\mathcal H}
\newcommand{\fM}{\mathfrak M}
\title{Anchored Likelihood-Ratio Geometry of Anonymous Shuffle Experiments:\\Exact Privacy Envelopes and Universal Low-Budget Design}
\author{Alex Shvets}
\date{Source: arXiv:2603.21197v1 (CC BY 4.0)}
\begin{document}
\maketitle

\noindent\emph{Source and licence.} Anchored Likelihood-Ratio Geometry of Anonymous Shuffle Experiments:\\Exact Privacy Envelopes and Universal Low-Budget Design. Version arXiv:2603.21197v1 (CC BY 4.0), distributed by arXiv under the Creative Commons Attribution 4.0 International licence, http://creativecommons.org/licenses/by/4.0/, as recorded in the arXiv licence metadata for this exact version and shown on the abstract page https://arxiv.org/abs/2603.21197v1. Full licence notice: \texttt{licenses/arxiv-2603.21197-CC-BY-4.0.txt}.\par\smallskip

\noindent\textit{Editorial scope.} Counted unit: the article's theorem thm:privacy-rigidity (source0.tex lines 859--877). The article's complete original proof of it is delivered with it (source0.tex lines 879--1005, one \texttt{proof} environment of 127 lines). No mathematics is written, repaired or re-derived: the statement and the proof are the article's own lines, in the article's own notation, and no retained line is substituted or re-flowed.  Retained uncounted as the source-local context the proof needs: lines 650--675.  Nothing was omitted from any retained span, so the package carries no omission record.  No retained line contains a citation, so the wrapper carries no bibliography and no bibliography entry.  No retained line contains a figure environment, an image-inclusion command or drawing code, so no image is delivered and the package declares no asset dependency.  Editorial formatting only, in addition: an AMS \texttt{amsart} preamble that restates the article's own declarations and macros used by the retained lines, namely the environments \texttt{proposition}, \texttt{theorem} on separate counters, together with the standard packages those lines need.  The retained environments print in this document as Proposition 1, Theorem 1 in order of appearance, whereas the article prints the same items with the numbering of its own full document.  The complete native source of the article as distributed in the arXiv e-print for this exact version is bundled unchanged as \texttt{sources/196923/source0.tex}.
\begin{proposition}[Pairwise likelihood ratios are scalar shadows]\label{prop:pairwise-lr}
Let $h,k\in\cH$, let $\rho\in\fM_d$, and assume $q_k\ll q_h$. Let $Y_h\sim\eta_h$. Then
\[
\frac{dq_k}{dq_h}(x)=\frac{1+k^\top x}{1+h^\top x}=1+(k-h)^\top T_h(x)
\qquad q_h\text{-a.s.}
\]
Hence, under $q_h$, the one-user likelihood ratio between $q_k$ and $q_h$ has law
\[
\mu_{h,k}:=\Law\bigl(1+(k-h)^\top Y_h\bigr).
\]
In particular, if $W_\rho$ is $\varepsilon_0$-LDP (so that all row pairs are mutually absolutely continuous), then for a row pair $i\neq j$,
\[
L_{ij}:=\frac{dW_\rho(\cdot\mid j)}{dW_\rho(\cdot\mid i)}(X)=1+(\gamma_j-\gamma_i)^\top Y_{\gamma_i},\qquad X\sim W_\rho(\cdot\mid i).
\]
If
\[
Q_{0,n}^{ij}:=W_\rho(\cdot\mid i)^{\otimes n},
\qquad
Q_{1,n}^{ij}:=\frac1n\sum_{m=1}^n W_\rho(\cdot\mid i)^{\otimes(m-1)}\otimes W_\rho(\cdot\mid j)\otimes W_\rho(\cdot\mid i)^{\otimes(n-m)},
\]
then under $Q_{0,n}^{ij}$,
\[
\frac{dQ_{1,n}^{ij}}{dQ_{0,n}^{ij}}=\frac1n\sum_{m=1}^n L_{ij,m},
\]
where $L_{ij,1},\dots,L_{ij,n}$ are i.i.d.\ with law $\mu_{ij}:=\Law(L_{ij})$.
\end{proposition}

\begin{theorem}[Privacy rigidity]\label{thm:privacy-rigidity}
Fix $\varepsilon_0\ge 0$ and $n\ge 1$. Let $\rho\in\fM_d$, let $i\neq j$, and assume that $W_\rho$ is $\varepsilon_0$-LDP. Suppose that for every $\alpha\ge 1$,
\[
H_\alpha\bigl(Q_{1,n}^{ij},Q_{0,n}^{ij}\bigr)
=
H_\alpha\bigl(Q_{1,n}^{\mathrm{BRR}},Q_{0,n}^{\mathrm{BRR}}\bigr)
\]
and
\[
H_\alpha\bigl(Q_{0,n}^{ij},Q_{1,n}^{ij}\bigr)
=
H_\alpha\bigl(Q_{0,n}^{\mathrm{BRR}},Q_{1,n}^{\mathrm{BRR}}\bigr).
\]
Then the one-user pairwise likelihood-ratio law satisfies
\[
\mu_{ij}=\mu^\star.
\]
Consequently, after merging outputs according to the two values of the pairwise likelihood ratio, the two-row reduction of $(W_\rho(\cdot\mid i),W_\rho(\cdot\mid j))$ is conditionally equivalent to binary randomized response with local parameter $\varepsilon_0$.
\end{theorem}

\begin{proof}
If $\varepsilon_0=0$, then the local constraint forces $W_\rho(\cdot\mid i)=W_\rho(\cdot\mid j)$, so $\mu_{ij}=\delta_1=\mu^\star$ and there is nothing to prove. Thus assume $\varepsilon_0>0$.

Let $\bar\mu_n$ be the law of
\[
\bar L_n:=\frac1n\sum_{m=1}^n L_{ij,m},
\]
where $L_{ij,1},\dots,L_{ij,n}$ are i.i.d.\ with law $\mu_{ij}$, and let $\bar\mu_n^\star$ be the corresponding averaged law built from $\mu^\star$. By \cref{prop:pairwise-lr},
\[
H_\alpha(Q_{1,n}^{ij},Q_{0,n}^{ij})=\int (t-\alpha)_+\,\bar\mu_n(dt),
\]
\[
H_\alpha(Q_{0,n}^{ij},Q_{1,n}^{ij})=\int (1-\alpha t)_+\,\bar\mu_n(dt),
\]
and likewise for $\bar\mu_n^\star$.

Define the stop-loss transform
\[
\mathsf S_{\bar\mu}(\alpha):=\int (t-\alpha)_+\,\bar\mu(dt),\qquad \alpha\in\R.
\]
Since both $\bar\mu_n$ and $\bar\mu_n^\star$ are compactly supported on $[e^{-\varepsilon_0},e^{\varepsilon_0}]$, their stop-loss transforms are convex and Lipschitz on $\R$. The right derivative of a convex function exists at every point, and for the stop-loss transform it satisfies
\[
\frac{d^+}{d\alpha}\mathsf S_{\bar\mu}(\alpha)=-\bar\mu((\alpha,\infty))
\qquad \forall\,\alpha\in\R.
\]
Therefore the equality
\[
\mathsf S_{\bar\mu_n}(\alpha)=\mathsf S_{\bar\mu_n^\star}(\alpha)\qquad \forall\,\alpha\ge 1
\]
implies
\[
\bar\mu_n((\alpha,\infty))=\bar\mu_n^\star((\alpha,\infty))
\qquad \forall\,\alpha\ge 1,
\]
This determines the common restriction of both measures to $(1,e^{\varepsilon_0}]$; the atom at $1$ is recovered only after combining this information with the lower stop-loss half below.

Now define the lower stop-loss transform
\[
\mathsf T_{\bar\mu}(\beta):=\int (\beta-t)_+\,\bar\mu(dt),\qquad \beta\in\R.
\]
It is also convex and Lipschitz, with right derivative
\[
\frac{d^+}{d\beta}\mathsf T_{\bar\mu}(\beta)=\bar\mu((-\infty,\beta])
\qquad \forall\,\beta\in\R.
\]
Since
\[
(1-\alpha t)_+=\alpha\Bigl(\frac1\alpha-t\Bigr)_+,
\]
the equality of the reverse hockey-stick profiles is equivalent to
\[
\mathsf T_{\bar\mu_n}(\beta)=\mathsf T_{\bar\mu_n^\star}(\beta)
\qquad\forall\,\beta\in[e^{-\varepsilon_0},1].
\]
Therefore
\[
\bar\mu_n((-\infty,\beta])=\bar\mu_n^\star((-\infty,\beta])
\qquad \forall\,\beta\in[e^{-\varepsilon_0},1],
\]
which determines both measures on $[e^{-\varepsilon_0},1]$ including any atoms. Hence
\[
\bar\mu_n=\bar\mu_n^\star.
\]

Let
\[
\phi(z):=\int e^{izt}\,\mu_{ij}(dt),
\qquad
\psi(z):=\int e^{izt}\,\mu^\star(dt)
\]
be the characteristic functions of the one-user laws. Since both laws are compactly supported, $\phi$ and $\psi$ extend to entire functions on $\mathbb C$. The characteristic function of $\bar\mu_n$ is
\[
\phi_{\bar\mu_n}(z)=\phi(z/n)^n,
\]
and similarly
\[
\phi_{\bar\mu_n^\star}(z)=\psi(z/n)^n.
\]
Because $\bar\mu_n=\bar\mu_n^\star$,
\[
\phi(z/n)^n=\psi(z/n)^n\qquad \forall\, z\in\R.
\]
Equivalently,
\[
\phi(w)^n=\psi(w)^n\qquad \forall\,w\in\R.
\]
The entire function
\[
F(w):=\phi(w)^n-\psi(w)^n
\]
vanishes on $\R$ (a set with limit points), hence vanishes identically on $\mathbb C$ by the identity theorem for entire functions.

Since $\psi(0)=1$, there exists a connected neighborhood $U$ of $0$ on which $\psi$ has no zeros. On $U$ the quotient
\[
r(w):=\frac{\phi(w)}{\psi(w)}
\]
is analytic and satisfies
\[
r(w)^n=1\qquad \forall\,w\in U.
\]
The image $r(U)$ is connected and contained in the finite set of $n$th roots of unity, so $r$ is constant on $U$. Because $r(0)=1$, we have $r\equiv 1$ on $U$. Hence $\phi=\psi$ on $U$, and therefore on all of $\mathbb C$ by the identity theorem. So
\[
\mu_{ij}=\mu^\star.
\]

Finally, since $\mu^\star$ is supported on the two points $e^{-\varepsilon_0}$ and $e^{\varepsilon_0}$, the pairwise likelihood ratio
\[
L_{ij}=\frac{dW_\rho(\cdot\mid j)}{dW_\rho(\cdot\mid i)}
\]
takes exactly these two values $W_\rho(\cdot\mid i)$-almost surely. Merge outputs according to the value of $L_{ij}$. Under row $i$, the resulting binary law has probabilities
\[
\frac{e^{\varepsilon_0}-1}{e^{\varepsilon_0}-e^{-\varepsilon_0}}
\quad\text{and}\quad
\frac{1-e^{-\varepsilon_0}}{e^{\varepsilon_0}-e^{-\varepsilon_0}}.
\]
Under row $j$, the probabilities are obtained by multiplying by the likelihood ratio:
\[
e^{-\varepsilon_0}\cdot \frac{e^{\varepsilon_0}-1}{e^{\varepsilon_0}-e^{-\varepsilon_0}}
=
\frac{1}{1+e^{\varepsilon_0}},
\qquad
e^{\varepsilon_0}\cdot \frac{1-e^{-\varepsilon_0}}{e^{\varepsilon_0}-e^{-\varepsilon_0}}
=
\frac{e^{\varepsilon_0}}{1+e^{\varepsilon_0}}.
\]
Thus the merged two-row experiment is exactly binary randomized response.
\end{proof}

\end{document}
